\chapter*{Resumen}
Robot Operating System 2 es un conjunto de bibliotecas de software y herramientas para construir soluciones de robótica.
Se basa en una arquitectura modular y conducida por eventos mediante el desarrollo de nodos 
que son interconectados mediante mecanismos de comunicación al tiempo de ejecución.

Los principios de diseño de ROS 2 son de especial beneficio para el diseño y escalado de sistemas complejos
donde multiples subsistemas conforman una solución y existen relaciones de dependencias no lineales entre ellos,
como es común de observar en el campo de la robotica, con entradas simultaneas desde varios sensores,
procesamiento en varios subsistemas tales como sistemas de visión por computador, planificación de caminos, cinematica inversa y control, 
multiples salidas desde actuaciones de sistemas electricos hasta comunicación con servidores u otros robots.

Sí bien hay una serie de ventajas al usar ROS 2, un aspecto poco explorado es el
impacto que tiene su uso en el rendimiento de aplicaciones. Al desarrollar una aplicación
modular con ROS 2 se esta haciendo uso implicito de un middleware de comunicaciónes llamado Data Distribution Service
o DDS, que se encarga de facilitar los mecanismos de comunicación interproceso, comparado a una aplicación uniproceso
los datos que un nodo envia a otro deben pasar por una pila completa de comunicación, lo cual alza 
preocupación sobre la degradación del rendimiento debido a un uso mayor de recursos computaciones para
implementar estos mecanismos, así como tambien una reducción en metricas clave como latencia y thoughput como consecuencia de ello.

En este trabajo se busca caracterizar el impacto sobre el rendimiento de una aplicación al usar ROS 2, 
para este fin se diseña e implementa una aplicación de visión por computador representativa de sistemas
usados en la industria, en una versión monolitica sin ROS 2 y una versión modular con ROS 2 sobre las cuales se 
realizaran pruebas de rendimiento.

Para las pruebas de rendimiento se prepara una metodologia de medición basada en herramientas de monitoreo de sistema e instrumentación de software,
y se monitorea el rendimiento de ambas versiones de la aplicación corriendo sobre una plataforma de hardware Nvidia Jetson y
posteriormente se analizan los resultados para evaluar el impacto del uso de ROS 2 sobre el rendimiento de la aplicación.

